% !TEX root = ../main.tex
\begin{ouverture}[synthese]{12}{Stratégies de synthèse organique}{Partie 1 · Constitution et transformations de la matière}
{Étapes d'une synthèse\enspace·\enspace Rendement\enspace·\enspace Substitution, addition, élimination\enspace·\enspace Protection et chimie verte}
\objectif{Identifier les étapes d'un protocole : transformation, isolement, purification, identification.}
\objectif{Justifier un montage à reflux et les choix expérimentaux d'une synthèse.}
\objectif{Déterminer le rendement d'une synthèse, y compris multiétapes.}
\objectif{Classer une réaction : substitution, addition, élimination ; modification de chaîne ou de groupe.}
\objectif{Comprendre une stratégie de protection et les principes d'une chimie plus respectueuse de l'environnement.}
\prerequis{Familles fonctionnelles et groupes caractéristiques (chapitre~11).}
\prerequis{Facteurs cinétiques et catalyse (chapitre~4) ; équilibre et estérification (chapitre~7).}
\prerequis{Quantité de matière, réactif limitant.}
\end{ouverture}
\essentiel
\essSection{1}{Les étapes d'une synthèse}
\begin{definition}[Les quatre étapes]
Une synthèse organique s'organise en quatre étapes : la \textbf{transformation} des réactifs (souvent sous chauffage à
reflux), l'\textbf{isolement} du produit brut (filtration, décantation, extraction), sa \textbf{purification}
(recristallisation, distillation, lavages) et son \textbf{identification} (CCM, température de fusion, spectroscopie
infrarouge), accompagnée si besoin d'un contrôle de pureté.
\end{definition}
\schema{12-etapes}
\begin{propriete}[Le chauffage à reflux]
Chauffer accélère la transformation (facteur cinétique) ; le réfrigérant vertical condense les vapeurs, qui retombent
dans le ballon : on chauffe ainsi \textbf{sans perte de matière}, même à ébullition.
\end{propriete}
\schema{12-montage}
\essSection{2}{Le rendement}
\begin{definition}[Rendement d'une synthèse]
Le rendement compare la quantité de produit réellement obtenue à la quantité maximale attendue, déterminée par le
réactif limitant :
\[\eta=\frac{n_{\text{obtenue}}}{n_{\max}}=\frac{m_{\text{obtenue}}}{m_{\max}},\qquad 0\leq\eta\leq1.\]
Une transformation non totale, des réactions parasites et les pertes de l'isolement et de la purification font que
$\eta<1$.
\end{definition}
\begin{exemple}
Pour synthétiser de l'aspirine (\ce{C9H8O4}, $M=\qty{180,2}{g.mol^{-1}}$), on part de 5,00~g d'acide salicylique
(\ce{C7H6O3}, $M=\qty{138,1}{g.mol^{-1}}$), réactif limitant : $n_{\max}=\frac{5{,}00}{138{,}1}=\qty{3,62e-2}{mol}$,
soit $m_{\max}=\num{3.62e-2}\times180{,}2=\qty{6,52}{g}$. On isole 4,2~g d'aspirine pure :
$\eta=\frac{4{,}2}{6{,}52}=0{,}64$, soit 64~\%.
\end{exemple}
\begin{codecote}[0.56]
\begin{pythoncol}
etapes = [0.85, 0.90, 0.70]
r = 1.0
for e in etapes:
    r = r * e
print(round(r, 3))
\end{pythoncol}
\begin{consolecol}
0.535
\end{consolecol}
\cote
\fontsize{9.5}{12.5}\selectfont
\capacite{Capacité numérique}\par\smallskip
Pour une synthèse en plusieurs étapes, le rendement global est le \textbf{produit} des rendements de chaque étape :
trois étapes pourtant honorables (85~\%, 90~\%, 70~\%) ne conservent que 54~\% du réactif de départ. Réduire le nombre
d'étapes est l'un des premiers leviers d'une synthèse efficace.
\end{codecote}
\essSection{3}{Catégories de réactions et stratégie}
\begin{definition}[Substitution, addition, élimination]
Dans une \textbf{substitution}, un atome ou un groupe est remplacé par un autre : \ce{R-Br + HO- -> R-OH + Br-}. Dans
une \textbf{addition}, deux réactifs se lient grâce à l'ouverture d'une liaison multiple :
\ce{CH2=CH2 + H2O -> CH3-CH2OH}. Dans une \textbf{élimination}, le substrat perd un petit groupe d'atomes et forme une
liaison multiple : \ce{CH3-CHOH-CH3 -> CH3-CH=CH2 + H2O}. Selon qu'elle touche le squelette ou le groupe
caractéristique, la réaction modifie la \textbf{chaîne} ou le \textbf{groupe}.
\end{definition}
\begin{propriete}[Protéger une fonction]
Lorsqu'un réactif porte plusieurs groupes réactifs, une réaction peut se produire au mauvais endroit. On
\textbf{protège} alors le groupe à préserver (on le transforme en un groupe inerte), on réalise la réaction souhaitée,
puis on \textbf{déprotège}. C'est la stratégie de la synthèse peptidique : pour lier l'acide aminé A à l'acide aminé B
dans l'ordre choisi, on bloque le groupe amine de A et le groupe carboxyle de B avant le couplage.
\end{propriete}
\schema{12-protection}
\essSection{4}{Optimiser et verdir une synthèse}
\begin{propriete}[Les leviers d'optimisation]
Pour accélérer : élever la température, utiliser un catalyseur (chapitre~4). Pour augmenter le rendement d'une
transformation non totale : introduire un réactif bon marché en excès, ou éliminer un produit au fur et à mesure, ce
qui maintient $Q_\mathrm{r}<K$ (chapitre~7). Un catalyseur, lui, ne modifie jamais l'état final.
\end{propriete}
\begin{definition}[Économie d'atomes]
L'économie d'atomes d'une réaction mesure la part de la masse des réactifs que l'on retrouve dans le produit visé :
\[\mathrm{EA}=\frac{M(\text{produit visé})}{\sum M(\text{réactifs})}\times100,\]
les masses molaires étant comptées avec les nombres stœchiométriques. Une addition a une économie d'atomes de
100~\% ; une substitution en a toujours une plus faible, car elle rejette un sous-produit.
\end{definition}
\begin{propriete}[Vers une chimie plus verte]
Une synthèse plus respectueuse de l'environnement cherche à : réduire le nombre d'étapes et les déchets (économie
d'atomes élevée), remplacer les réactifs et solvants dangereux, travailler à température modérée, catalyser, et
utiliser des matières premières renouvelables.
\end{propriete}
\begin{rediger}
Pour un rendement, citer le réactif limitant, la quantité maximale attendue puis la quantité obtenue, et conclure en
pourcentage. Pour une catégorie de réaction, nommer la liaison créée ou rompue : « la double liaison C$=$C s'ouvre :
c'est une addition ».
\end{rediger}
\exercices

\rubrique{Découverte}
\begin{blocLivre}{ex-12-1}
\begin{exercice}[1]{Les étapes d'un protocole}
Protocole de synthèse de l'éthanoate de 3-méthylbutyle (arôme de banane) : « \textit{(1) Dans un ballon, introduire
20~mL de 3-méthylbutan-1-ol, 30~mL d'acide éthanoïque, quelques gouttes d'acide sulfurique et de la pierre ponce.
(2) Chauffer à reflux 40~minutes. (3) Verser le contenu refroidi dans de l'eau salée glacée, puis séparer la phase
organique à l'ampoule à décanter. (4) La laver avec une solution d'hydrogénocarbonate de sodium, puis la sécher.
(5) Vérifier la pureté par chromatographie sur couche mince.} »
\begin{questions}
\item Attribuer chacune des phrases numérotées à l'une des quatre étapes d'une synthèse.
\item Quel est le rôle de l'acide sulfurique ? celui de la pierre ponce ?
\item Pourquoi chauffe-t-on à reflux plutôt que dans un bécher ouvert ?
\item À quelle famille le produit appartient-il ?
\end{questions}
\end{exercice}
\end{blocLivre}
\begin{blocLivre}{ex-12-2}
\begin{exercice}[1]{Le montage à reflux}
Le schéma représente un montage de chauffage à reflux dont les légendes ont été remplacées par des lettres.
\schema{12-ex02}
\begin{questions}
\item Nommer les éléments A à E.
\item Dans quel sens l'eau circule-t-elle dans l'élément B ? Pourquoi ?
\item Expliquer la phrase : « le reflux permet de chauffer sans perdre de matière ».
\item Quel élément du montage joue un rôle de sécurité vis-à-vis des vapeurs inflammables ?
\end{questions}
\end{exercice}
\end{blocLivre}
\begin{blocLivre}{ex-12-3}
\begin{exercice}[1]{Calculer un rendement}
Lors d'une synthèse de l'éthanoate d'éthyle \ce{CH3-COO-C2H5}, le réactif limitant est l'acide éthanoïque, introduit à
raison de 0,100~mol. Après purification, on isole 5,9~g d'ester.
\begin{donnees}
$M(\text{ester})=\qty{88,1}{g.mol^{-1}}$.
\end{donnees}
\begin{questions}
\item Écrire l'équation de la réaction d'estérification entre l'acide éthanoïque et l'éthanol.
\item Quelle est la quantité maximale d'ester attendue ? la masse correspondante ?
\item Calculer le rendement de la synthèse.
\item Citer deux raisons pour lesquelles le rendement n'atteint pas 100~\%.
\end{questions}
\end{exercice}
\end{blocLivre}
\begin{blocLivre}{ex-12-4}
\begin{exercice}[1]{Substitution, addition ou élimination ?}
On considère les transformations suivantes.
\begin{questions}[label=\otLabel{\Alph*.},leftmargin=1.9em,itemsep=0pt,parsep=0.15em]
\item \ce{CH2=CH2 + H2O -> CH3-CH2-OH}
\item \ce{CH3-CH2-CH2-Br + HO- -> CH3-CH2-CH2-OH + Br-}
\item \ce{CH3-CHOH-CH3 -> CH3-CH=CH2 + H2O}
\item \ce{CH2=CH2 + H2 -> CH3-CH3}
\end{questions}
\begin{questions}
\item Classer chaque transformation : substitution, addition ou élimination. Justifier.
\item Pour la transformation B, préciser l'atome ou le groupe remplacé et le groupe entrant.
\item Laquelle de ces transformations est mise en jeu dans une polymérisation ?
\end{questions}
\end{exercice}
\end{blocLivre}
\begin{blocLivre}{ex-12-5}
\begin{exercice}[1]{Chaîne ou groupe ?}
Pour chaque transformation, indiquer si elle modifie la chaîne carbonée, le groupe caractéristique, ou les deux.
\begin{questions}[label=\otLabel{\alph*.},leftmargin=1.9em,itemsep=0pt,parsep=0.15em]
\item \ce{CH3-CH2-OH -> CH3-COOH} (fabrication du vinaigre) ;
\item $n\,\ce{CH2=CH2 ->}\ \ce{-(CH2-CH2)_{$n$}-}$ (polymérisation) ;
\item \ce{CH3-CH2-CH2-OH -> CH3-CH2-CH2-Br} ;
\item \ce{CH3-COOH + CH3-OH -> CH3-COO-CH3 + H2O}.
\end{questions}
\begin{questions}
\item Justifier chaque réponse en comparant squelettes et groupes caractéristiques.
\item Parmi ces transformations, laquelle est une substitution ?
\end{questions}
\end{exercice}
\end{blocLivre}
\rubrique{Entraînement}
\begin{blocLivre}{ex-12-6}
\begin{exercice}[2]{La synthèse de l'aspirine}
Au laboratoire, on synthétise l'aspirine \ce{C9H8O4} à partir de 5,00~g d'acide salicylique \ce{C7H6O3} et d'un excès
d'anhydride éthanoïque, sous chauffage à reflux pendant 20~minutes. Le mélange refroidi est versé sur de la glace :
l'aspirine cristallise. Après filtration sur Büchner, recristallisation et séchage, on pèse 4,2~g de cristaux, dont la
température de fusion mesurée vaut 135~°C (valeur attendue : 136~°C).
\begin{donnees}
$M(\ce{C7H6O3})=\qty{138,1}{g.mol^{-1}}$ ; $M(\ce{C9H8O4})=\qty{180,2}{g.mol^{-1}}$ ; une mole d'acide salicylique
donne au plus une mole d'aspirine.
\end{donnees}
\begin{questions}
\item Identifier, dans le protocole, les quatre étapes de la synthèse.
\item Pourquoi l'anhydride éthanoïque est-il introduit en excès ?
\item Calculer la masse maximale d'aspirine attendue, puis le rendement.
\item Que permet de conclure la mesure de la température de fusion ?
\end{questions}
\end{exercice}
\end{blocLivre}
\begin{blocLivre}{ex-12-7}
\begin{exercice}[2]{Choisir les bonnes conditions}
On étudie l'estérification \ce{CH3COOH + C2H5OH <=> CH3COOC2H5 + H2O}, lente et non totale ($K=4{,}0$, chapitre~7).
Quatre propositions sont faites pour améliorer la synthèse :
\begin{questions}[label=\otLabel{\alph*.},leftmargin=1.9em,itemsep=0pt,parsep=0.15em]
\item ajouter quelques gouttes d'acide sulfurique ;
\item chauffer le mélange ;
\item introduire l'éthanol en large excès ;
\item éliminer l'eau au fur et à mesure de sa formation.
\end{questions}
\begin{questions}
\item Pour chaque proposition, indiquer si elle augmente la vitesse de la transformation, son rendement final, ou les
deux. Justifier à l'aide des chapitres 4 et 7.
\item Pourquoi un catalyseur ne peut-il pas augmenter le rendement final ?
\item Quelle grandeur du chapitre 7 l'élimination de l'eau fait-elle varier, et dans quel sens ?
\end{questions}
\end{exercice}
\end{blocLivre}
\begin{blocLivre}{ex-12-8}
\begin{exercice}[2]{L'économie d'atomes}
Deux voies de synthèse du propan-2-ol sont envisagées :
\begin{questions}[label=\otLabel{\Alph*.},leftmargin=1.9em,itemsep=0pt,parsep=0.15em]
\item addition d'eau sur le propène : \ce{CH3-CH=CH2 + H2O -> CH3-CHOH-CH3} ;
\item substitution sur le 2-bromopropane : \ce{CH3-CHBr-CH3 + NaOH -> CH3-CHOH-CH3 + NaBr}.
\end{questions}
\begin{donnees}
$M(\ce{C3H6})=\qty{42,1}{g.mol^{-1}}$ ; $M(\ce{H2O})=\qty{18,0}{g.mol^{-1}}$ ;
$M(\ce{C3H7Br})=\qty{123,0}{g.mol^{-1}}$ ; $M(\ce{NaOH})=\qty{40,0}{g.mol^{-1}}$ ;
$M(\ce{C3H8O})=\qty{60,1}{g.mol^{-1}}$.
\end{donnees}
\begin{questions}
\item Rappeler la définition de l'économie d'atomes.
\item Calculer l'économie d'atomes de chaque voie.
\item Quel sous-produit la voie B rejette-t-elle ? Commenter.
\item Pourquoi une addition a-t-elle toujours une économie d'atomes de 100~\% ?
\end{questions}
\end{exercice}
\end{blocLivre}
\begin{blocLivre}{ex-12-9}
\begin{exercice}[2]{Le rendement d'une synthèse multiétapes}
Une molécule est préparée en trois étapes successives, de rendements respectifs 85~\%, 90~\% et 70~\%. Le programme
calcule le rendement global et la quantité finale obtenue à partir de 0,50~mol de réactif.
\begin{python}
etapes = [0.85, 0.90, 0.70]
r = 1.0
for e in etapes:
    r = r * e
print(round(r, 3), round(0.50 * r, 3))
\end{python}
\begin{console}
0.535 0.268
\end{console}
\begin{questions}
\item \capacite{Capacité numérique}Justifier que le rendement global est le produit des rendements des étapes.
\item Interpréter les deux nombres affichés.
\item L'équipe peut porter une seule étape à 95~\%. Laquelle choisir pour le meilleur gain ? Calculer le nouveau
rendement global.
\item Quel serait le rendement global d'une synthèse équivalente en une seule étape de rendement 70~\% ? Conclure.
\end{questions}
\end{exercice}
\end{blocLivre}
\begin{blocLivre}{ex-12-10}
\begin{exercice}[2]{Pourquoi protéger ?}
On veut synthétiser le dipeptide Gly--Ala, dans lequel le groupe carboxyle de la glycine (Gly) est lié au groupe amine
de l'alanine (Ala). Or chaque acide aminé porte à la fois un groupe amine et un groupe carboxyle.
\schema{12-ex10}
\begin{questions}
\item Expliquer pourquoi un simple mélange de Gly et d'Ala conduit à quatre dipeptides différents. Les désigner.
\item Quelle stratégie permet d'obtenir uniquement Gly--Ala ? Décrire les étapes dans l'ordre.
\item Pourquoi faut-il, à la fin, une étape de déprotection ?
\item Chaque étape supplémentaire a un rendement imparfait. Quel est le prix à payer de cette stratégie ?
\end{questions}
\end{exercice}
\end{blocLivre}
\begin{blocLivre}{ex-12-11}
\begin{exercice}[2]{La juste durée du reflux}
Pour une estérification conduite à reflux, on mesure le rendement obtenu en arrêtant la synthèse à différentes durées
de chauffage.
\schema{12-ex11}
\begin{questions}
\item Décrire l'évolution du rendement avec la durée de chauffage. Pourquoi apparaît-il un palier ?
\item Quelle durée de chauffage recommander ? Justifier.
\item Reproduire l'allure de la courbe et tracer celle attendue avec un catalyseur. Justifier.
\item Quel levier faudrait-il utiliser pour dépasser la valeur du palier ?
\end{questions}
\end{exercice}
\end{blocLivre}
\rubrique{Objectif bac}
\begin{blocLivre}{ex-12-12}
\begin{exercice}[3]{Deux routes vers l'ibuprofène}
\begin{docu}{1}{Des années 1960 aux années 1990}
L'ibuprofène a longtemps été produit par la synthèse « Boots » : six étapes, et une économie d'atomes globale de
40~\% seulement. Dans les années 1990, le procédé « BHC » l'a remplacée : trois étapes catalytiques, dont les
catalyseurs sont récupérés et réutilisés, et une économie d'atomes de 77~\%. Ce procédé a reçu un prix de chimie
verte.
\end{docu}
\begin{donnees}
$M(\text{ibuprofène})=\qty{206,3}{g.mol^{-1}}$ ; somme des masses molaires des réactifs engagés :
\qty{514,5}{g.mol^{-1}} pour la voie Boots, \qty{266,8}{g.mol^{-1}} pour la voie BHC.
\end{donnees}
\begin{questions}
\item Vérifier les deux valeurs d'économie d'atomes annoncées.
\item \initiative Pour une tonne d'ibuprofène produite, estimer la masse de sous-produits rejetés par chaque voie.
\item En quoi la réduction du nombre d'étapes améliore-t-elle aussi le rendement global ?
\item Citer deux principes de la chimie verte illustrés par le procédé BHC.
\end{questions}
\end{exercice}
\end{blocLivre}
\begin{blocLivre}{ex-12-13}
\begin{exercice}[3]{Assembler un dipeptide sans fausse note}
\begin{docu}{1}{Le plan de synthèse}
Pour préparer le dipeptide Ala--Gly, un laboratoire suit le plan suivant, chaque étape ayant un rendement de 90~\% :
(1) protection du groupe amine de l'alanine ; (2) protection du groupe carboxyle de la glycine ; (3) couplage : le
groupe carboxyle libre de l'alanine réagit avec le groupe amine libre de la glycine ; (4) déprotections, qui libèrent
le dipeptide.
\end{docu}
\begin{questions}
\item Quel groupe caractéristique le couplage crée-t-il ? De quelle famille, vue au chapitre~11, le dipeptide
relève-t-il ?
\item Pourquoi les étapes 1 et 2 sont-elles indispensables ? Que se formerait-il sans elles ?
\item Justifier l'ordre des protections choisi : pourquoi bloque-t-on l'amine de l'alanine et non son carboxyle ?
\item \initiative Calculer le rendement global du plan de synthèse.
\item La synthèse industrielle des peptides enchaîne des dizaines de couplages. Commenter la contrainte que fait peser
le rendement de chaque étape.
\end{questions}
\end{exercice}
\end{blocLivre}
\begin{blocLivre}{ex-12-14}
\begin{exercice}[3]{La synthèse du paracétamol}
\begin{docu}{1}{Le protocole}
Dans un ballon, dissoudre 5,45~g de 4-aminophénol dans 15~mL d'eau tiède ; ajouter lentement 6,0~mL d'anhydride
éthanoïque (excès), puis chauffer à reflux 15~minutes. Verser dans l'eau glacée : le paracétamol cristallise. Filtrer
sur Büchner, recristalliser dans l'eau chaude, sécher à l'étuve. Masse obtenue : 6,0~g. Température de fusion
mesurée : 169~°C (valeur de référence : 170~°C).
\end{docu}
\begin{donnees}
$M(\text{4-aminophénol})=\qty{109,1}{g.mol^{-1}}$ ; $M(\text{paracétamol})=\qty{151,2}{g.mol^{-1}}$ ; une mole de
4-aminophénol donne au plus une mole de paracétamol.
\end{donnees}
\begin{questions}
\item Le 4-aminophénol porte un groupe hydroxyle et un groupe amine ; l'anhydride réagit ici avec le groupe amine pour
former un amide. Quel groupe du paracétamol en résulte ?
\item Identifier les quatre étapes de la synthèse dans le protocole.
\item Calculer la masse maximale de paracétamol attendue, puis le rendement.
\item \initiative La recristallisation fait perdre environ 10~\% du produit brut. Quel rendement aurait-on sans
elle ? Pourquoi la conserve-t-on malgré tout ?
\end{questions}
\end{exercice}
\end{blocLivre}
\begin{blocLivre}{ex-12-15}
\begin{exercice}[3]{Une estérification plus verte}
\begin{docu}{1}{Deux procédés pour un même ester}
Procédé A, classique : acide éthanoïque et éthanol, acide sulfurique concentré, chauffage à reflux à 80~°C pendant
une heure ; rendement 67~\%. Procédé B, enzymatique : les mêmes réactifs, une lipase immobilisée comme catalyseur,
40~°C pendant trois heures, sans autre solvant ; rendement 90~\% grâce à l'élimination de l'eau par un tamis
moléculaire ; l'enzyme est réutilisée plus de vingt fois.
\end{docu}
\begin{donnees}
$M(\ce{CH3COOH})=\qty{60,1}{g.mol^{-1}}$ ; $M(\ce{C2H5OH})=\qty{46,1}{g.mol^{-1}}$ ;
$M(\ce{CH3COOC2H5})=\qty{88,1}{g.mol^{-1}}$.
\end{donnees}
\begin{questions}
\item Écrire l'équation de la réaction et calculer son économie d'atomes. Dépend-elle du procédé ?
\item Expliquer pourquoi le procédé B atteint 90~\% de rendement malgré une constante d'équilibre inchangée.
\item Comparer les deux procédés du point de vue de l'énergie et de la sécurité.
\item \initiative Relever, dans le procédé B, trois choix qui relèvent de la chimie verte.
\end{questions}
\end{exercice}
\end{blocLivre}
